\begin{activite}[Produit d'un nombre négatif par un nombre positif] \label{MultDivRelatifs_acti1}

On considère l'expression $A = (-2) + (-2) + (-2) + (-2)$.

\begin{partie}
Quelle est la valeur de $A$ ? \\[0.5em]
On va revenir sur le sens de la multiplication : $20 + 20 + 20$ est la somme de trois termes tous égaux. On peut donc écrire cette somme sous la forme du produit $20 \times 3$ qui se lit « 20 multiplié par 3 ».
\end{partie}

\begin{partie}
Écris $A$ sous la forme d'un produit.
 \end{partie}
 
\begin{partie} 
Écris les expressions suivantes sous la forme d'une somme et calcule-les :
 \begin{colenumerate}{4}
  \item $(-6) \times 3$ ;
  \item $(-22) \times 5$ ;
  \item $(-7) \times 7$ ;
  \item $(-1,5) \times 6$.
  \end{colenumerate}
 \end{partie}
 
\begin{partie} \label{MultDivRelatifs_acti2}
Trouve une règle permettant de calculer le produit d'un nombre négatif par un nombre positif.
 \end{partie}

\end{activite}

%%%%%%%%%%%%%%%%%%%%%%%%%%%%%%%%%%%%%%%%%%%%%%%%%%%%%%%%%%%%%%%%%%%%%%%%%

\begin{activite}[À propos des produits]

\begin{minipage}[c]{0.66\linewidth}
\begin{partie}
Voici une table de multiplication :
\begin{enumerate}
 \item Recopie-la sur ton cahier et complète la partie qui concerne le produit de deux nombres positifs (en bas à droite).
 \item D'après le résultat de la partie \ref{MultDivRelatifs_acti2} de l'activité \ref{MultDivRelatifs_acti1}, complète la partie qui concerne le produit d'un nombre négatif par un nombre positif (en haut à droite).
 \item Observe les résultats dans cette table de multiplication et complète-la entièrement, en expliquant tes choix.
 \item À l'aide d'un tableur, crée cette table de multiplication et vérifie que les résultats obtenus sont les mêmes que les tiens.
 \end{enumerate}
\end{partie}
\end{minipage} \hfill%
 \begin{minipage}[c]{0.3\linewidth}
  \includegraphics[width=5cm]{grille_produits}
  \end{minipage} \\

\begin{partie}
Application sur quelques exemples :
\begin{enumerate}
 \item En t'aidant de la table, donne le résultat pour chaque calcul suivant :
 \begin{colitemize}{4}
  \item $A = (-5) \times 4$ ;
  \item $B = 3 \times (-2)$ ;
  \item $C = 5 \times (-4)$ ;
  \item $D = (-1) \times (-3)$.
  \end{colitemize}
 \item En t'inspirant de ce qui précède, propose un résultat pour les calculs suivants :
 \begin{colitemize}{4}
  \item $E = (-9,2) \times 2$ ;
  \item $F = 1,5 \times (-8)$ ;
  \item $G = (-3,14) \times 0$ ;
  \item $H = (-1,2) \times (-0,1)$.
  \end{colitemize}
 \item Vérifie ces résultats à la calculatrice.
 \end{enumerate}
\end{partie}

\begin{partie}
Propose une règle qui permet, dans tous les cas, de calculer le produit de deux nombres relatifs.
\end{partie}

\end{activite}

%%%%%%%%%%%%%%%%%%%%%%%%%%%%%%%%%%%%%%%%%%%%%%%%%%%%%%%%%%%%%%%%%%%%%%%%%

\begin{activite}[Produit de plusieurs nombres relatifs]

\begin{partie}
Calcule ces expressions et déduis-en une règle pour trouver rapidement chaque résultat :
\begin{itemize}
 \item $A = (-1) \times (-1)$ ;
 \item $B = (-1) \times (-1) \times (-1)$ ;
 \item $C = (-1) \times (-1) \times (-1) \times (-1)$ ;
 \item $D = (-1) \times (-1) \times (-1) \times (-1) \times (-1)$ ;
 \item $E = (-1) \times (-1) \times (-1) \times (-1) \times (-1) \times (-1)$.
 \end{itemize}
\end{partie}

\begin{partie}
On sait que $(-4) = (-1) \times 4$ et $(-2) = (-1) \times 2$.
\begin{enumerate}
 \item Complète alors le calcul suivant :
\begin{center} $(-4) \times (-2) \times (-5) = (-1) \times \ldots \times (-1) \times \ldots \times (- ) \times \ldots$ \end{center}
\begin{center} $\phantom{(-4) \times (-2) \times (-5) }= (-1) \times (-1) \times (-1) \times \ldots \times \ldots \times \ldots$ \end{center}
 \item Déduis-en une méthode pour trouver le résultat de $(-4) \times (-2) \times (-5)$.
 \end{enumerate}
\end{partie}

\begin{partie}
Inspire-toi de la question précédente pour effectuer le calcul suivant :
\begin{itemize}
 \item $F = (-2) \times (-3) \times 5 \times (-4) \times 6 \times (-5)$.
 \end{itemize}
\end{partie}

\begin{partie}
Propose une méthode pour multiplier plusieurs nombres relatifs.
\end{partie}

\end{activite}

%%%%%%%%%%%%%%%%%%%%%%%%%%%%%%%%%%%%%%%%%%%%%%%%%%%%%%%%%%%%%%%%%%%%%%%%%

\begin{activite}[Quotient de nombres relatifs]

Revenons sur le sens de la division : \\[0.5em]
Écrire $3 \times 5 = 15$ revient à écrire $3 = 15 : 5$ ou $5 = 15 : 3$.
       
\begin{partie}
Recopie et complète les trous par les nombres manquants pour que les égalités soient correctes :
\begin{colenumerate}{4}
 \item $4 \times  \ldots = 12$ ;
 \item $(-5) \times  \ldots = 130$ ;
 \item $8 \times  \ldots = (-16)$ ;
 \item $ \ldots \times (-3) = (-27)$.
 \end{colenumerate}
\end{partie}

\begin{partie}
Écris ces nombres manquants sous forme de quotients.
\end{partie}

\begin{partie}
Que dire du quotient de deux nombres relatifs ?
\end{partie}

\end{activite}

%%%%%%%%%%%%%%%%%%%%%%%%%%%%%%%%%%%%%%%%%%%%%%%%%%%%%%%%%%%%%%%%%%%%%%%%%
